% =============================================================================
% =============================================================================
%
%              EBF CHAPTER TEMPLATE: Comprehensive & Adaptive
%
% =============================================================================
% =============================================================================
%
% Version: 2.0 (January 2026)
% Status: Authoritative Template for All Main Document Chapters
%
% This template defines the COMPLETE structure for EBF chapters.
% All elements are documented; adapt based on chapter type.
%
% =============================================================================

% #############################################################################
%
%                    DOCUMENT ARCHITECTURE OVERVIEW
%
% #############################################################################
%
% The EBF document follows a dual structure:
%
% ┌─────────────────────────────────────────────────────────────────────────┐
% │                        MAIN DOCUMENT (Chapters)                         │
% │                                                                         │
% │  Purpose: Conceptual narrative, intuition, motivation                   │
% │  Style:   Accessible prose, minimal formalism                           │
% │  Length:  15-50 pages per chapter                                       │
% │  Reader:  Economists, policy makers, PhD students                       │
% └─────────────────────────────────────────────────────────────────────────┘
%                              │
%                              ▼
% ┌─────────────────────────────────────────────────────────────────────────┐
% │                         APPENDICES (50 total)                           │
% │                                                                         │
% │  Purpose: Formal proofs, complete axioms, full specifications          │
% │  Style:   Technical, mathematical, exhaustive                           │
% │  Length:  3-200 pages (varies)                                          │
% │  Reader:  Specialists, theorists, advanced researchers                  │
% └─────────────────────────────────────────────────────────────────────────┘
%
% =============================================================================


% #############################################################################
%
%                    CHAPTER TYPES & ADAPTATION GUIDE
%
% #############################################################################
%
% EBF has THREE chapter types. Adapt template accordingly:
%
% ┌─────────────────────────────────────────────────────────────────────────┐
% │  TYPE A: CORE Chapters (Ch. 5, 9, 10, 11, 12, 13)                       │
% │  ─────────────────────────────────────────────────────────────────────  │
% │  • Has dedicated CORE Appendix (B, V, C, AU, AV, AW)                    │
% │  • USE: CORE Connection Box, full formal sections                       │
% │  • USE: 9C Integration Table                                            │
% └─────────────────────────────────────────────────────────────────────────┘
%
% ┌─────────────────────────────────────────────────────────────────────────┐
% │  TYPE B: Foundation Chapters (Ch. 1-4, 6-8)                             │
% │  ─────────────────────────────────────────────────────────────────────  │
% │  • Builds conceptual groundwork                                         │
% │  • SKIP: CORE Connection Box (no dedicated CORE)                        │
% │  • USE: Standard subsection structure                                   │
% └─────────────────────────────────────────────────────────────────────────┘
%
% ┌─────────────────────────────────────────────────────────────────────────┐
% │  TYPE C: Application Chapters (Ch. 14-19)                               │
% │  ─────────────────────────────────────────────────────────────────────  │
% │  • Applies framework to practice                                        │
% │  • USE: Multiple worked examples                                        │
% │  • USE: Policy implications section                                     │
% └─────────────────────────────────────────────────────────────────────────┘
%
% =============================================================================


% #############################################################################
%
%                    COMPLETE CHAPTER STRUCTURE
%
% #############################################################################
%
%    Every EBF Chapter follows this master structure:
%
%    ┌─────────────────────────────────────────────────────────────────────┐
%    │  HEADER                                                             │
%    ├─────────────────────────────────────────────────────────────────────┤
%    │  □ Metadata comment block (version, purpose, appendices)            │
%    │  □ \chapter{} command with \label{}                                 │
%    │  □ Navigation Box (\BCMQuestionNav{})                               │
%    │  □ CORE Connection Box (for TYPE A chapters)                        │
%    │  □ Appendix References Box                                          │
%    ├─────────────────────────────────────────────────────────────────────┤
%    │  OPENING                                                            │
%    ├─────────────────────────────────────────────────────────────────────┤
%    │  □ Opening paragraph (2-3 sentences: what and why)                  │
%    │  □ Intuition box (motivating example)                               │
%    │  □ Central Question box                                             │
%    │  □ Chapter Overview (section roadmap)                               │
%    ├─────────────────────────────────────────────────────────────────────┤
%    │  CORE CONTENT (Multiple Sections)                                   │
%    ├─────────────────────────────────────────────────────────────────────┤
%    │  □ Subsections with logical flow                                    │
%    │  □ Definition environments for formal concepts                      │
%    │  □ Intuition environments for explanations                          │
%    │  □ Cross-reference environments for links                           │
%    │  □ Key equations (numbered, interpreted)                            │
%    │  □ Tables for structured information                                │
%    ├─────────────────────────────────────────────────────────────────────┤
%    │  WORKED EXAMPLES                                                    │
%    ├─────────────────────────────────────────────────────────────────────┤
%    │  □ At least one concrete example with numbers                       │
%    │  □ Before/after comparison if applicable                            │
%    │  □ Multi-level example if applicable                                │
%    ├─────────────────────────────────────────────────────────────────────┤
%    │  CLOSING                                                            │
%    ├─────────────────────────────────────────────────────────────────────┤
%    │  □ Summary (enumerated key points)                                  │
%    │  □ Formal Details paragraph (appendix references)                   │
%    │  □ What Comes Next paragraph (transition)                           │
%    │  □ Reading Path cross-reference box                                 │
%    └─────────────────────────────────────────────────────────────────────┘
%
% =============================================================================


% *****************************************************************************
% *****************************************************************************
%
%                    BEGIN TEMPLATE CODE
%
% *****************************************************************************
% *****************************************************************************


%%%%%%%%%%%%%%%%%%%%%%%%%%%%%%%%%%%%%%%%%%%%%%%%%%%%%%%%%%%%%%%%%%%%%%%%%%%%%%%
% SECTION 1: HEADER METADATA BLOCK
%%%%%%%%%%%%%%%%%%%%%%%%%%%%%%%%%%%%%%%%%%%%%%%%%%%%%%%%%%%%%%%%%%%%%%%%%%%%%%%
%
% REQUIRED for all chapters. Adapt fields as needed.
%
%------------------------------------------------------------------------------
% Chapter [NUMBER]: [Title]
%------------------------------------------------------------------------------
% Version: [X.Y] ([Month Year])
%
% Purpose: [2-3 sentence description of chapter's role in EBF]
%
% Primary Appendix: [CODE] - [CATEGORY-NAME] (e.g., AW - CORE-STAGE)
% Secondary Appendices: [List of related appendices]
%
% Prerequisites: [Chapters that should be read first]
% Leads to: [Chapters that build on this]
%
% Chapter Type: [A: CORE / B: Foundation / C: Application]
%------------------------------------------------------------------------------


%%%%%%%%%%%%%%%%%%%%%%%%%%%%%%%%%%%%%%%%%%%%%%%%%%%%%%%%%%%%%%%%%%%%%%%%%%%%%%%
% SECTION 2: CHAPTER COMMAND & LABEL
%%%%%%%%%%%%%%%%%%%%%%%%%%%%%%%%%%%%%%%%%%%%%%%%%%%%%%%%%%%%%%%%%%%%%%%%%%%%%%%

\chapter{[Chapter Title]}
\label{ch:[short-label]}


%%%%%%%%%%%%%%%%%%%%%%%%%%%%%%%%%%%%%%%%%%%%%%%%%%%%%%%%%%%%%%%%%%%%%%%%%%%%%%%
% SECTION 3: NAVIGATION BOX
%%%%%%%%%%%%%%%%%%%%%%%%%%%%%%%%%%%%%%%%%%%%%%%%%%%%%%%%%%%%%%%%%%%%%%%%%%%%%%%
%
% REQUIRED for all chapters. Shows position in the 9 BCM Questions.
%
%==============================================================================
% DIE 9 FRAGEN NAVIGATION BOX
%==============================================================================
\BCMQuestionNav{[CHAPTER-NUMBER]}


%%%%%%%%%%%%%%%%%%%%%%%%%%%%%%%%%%%%%%%%%%%%%%%%%%%%%%%%%%%%%%%%%%%%%%%%%%%%%%%
% SECTION 4: CORE CONNECTION BOX (TYPE A CHAPTERS ONLY)
%%%%%%%%%%%%%%%%%%%%%%%%%%%%%%%%%%%%%%%%%%%%%%%%%%%%%%%%%%%%%%%%%%%%%%%%%%%%%%%
%
% USE FOR: Chapters with dedicated CORE appendix (5, 9, 10, 11, 12, 13)
% SKIP FOR: Foundation chapters (1-4, 6-8) and Application chapters (14-19)
%
% COLOR CODING:
%   - CORE-HOW (B):    blue
%   - CORE-WHEN (V):   green
%   - CORE-WHAT (C):   red
%   - CORE-WHO (AAA):  purple
%   - CORE-AWARE (AU): orange
%   - CORE-READY (AV): violet
%   - CORE-STAGE (AW): teal
%
%==============================================================================
% CORE CONNECTION BOX
%==============================================================================
\begin{tcolorbox}[
    colback=[COLOR]!5!white,
    colframe=[COLOR]!75!black,
    title={\textbf{9C CORE Connection: [CORE-NAME] (Question [N])}},
    fonttitle=\bfseries
]
\textbf{The [N]th CORE Question:} [Question text]

\medskip
\begin{tabular}{ll}
\textbf{CORE Appendix:} & [CODE] ([CATEGORY-NAME]) \\
\textbf{Output:} & [What this CORE produces] \\
\textbf{Key Insight:} & [One-sentence summary] \\
\textbf{Novel Contribution:} & [What's new in EBF]
\end{tabular}

\medskip
\textit{[Brief statement of chapter's role in 9C framework]}
\end{tcolorbox}


%%%%%%%%%%%%%%%%%%%%%%%%%%%%%%%%%%%%%%%%%%%%%%%%%%%%%%%%%%%%%%%%%%%%%%%%%%%%%%%
% SECTION 5: APPENDIX REFERENCES BOX
%%%%%%%%%%%%%%%%%%%%%%%%%%%%%%%%%%%%%%%%%%%%%%%%%%%%%%%%%%%%%%%%%%%%%%%%%%%%%%%
%
% REQUIRED for all chapters. Links sections to relevant appendices.
%
%==============================================================================
% APPENDIX REFERENCES BOX
%==============================================================================
\begin{tcolorbox}[
    colback=yellow!10!white,
    colframe=orange!75!black,
    title={\textbf{Appendix References for This Chapter}},
    fonttitle=\bfseries
]
\begin{tabular}{lll}
\textbf{Appendix} & \textbf{Content} & \textbf{Section} \\
\midrule
[CODE] ([NAME]) & [Brief description] & \ref{sec:[label]} \\
[CODE] ([NAME]) & [Brief description] & \ref{sec:[label]} \\
F (REF-EXAMPLES) & Worked examples & \ref{sec:[label]-applications} \\
G (REF-GLOSSARY) & Symbol definitions & All sections
\end{tabular}
\end{tcolorbox}


%%%%%%%%%%%%%%%%%%%%%%%%%%%%%%%%%%%%%%%%%%%%%%%%%%%%%%%%%%%%%%%%%%%%%%%%%%%%%%%
%                          PART I: OPENING
%%%%%%%%%%%%%%%%%%%%%%%%%%%%%%%%%%%%%%%%%%%%%%%%%%%%%%%%%%%%%%%%%%%%%%%%%%%%%%%

\section{Introduction: [Subtitle]}
\label{sec:[label]-intro}

%------------------------------------------------------------------------------

% Opening paragraph: WHAT this chapter does and WHY it matters
% Keep to 2-3 sentences. Accessible language.

[Opening paragraph that introduces the core question of this chapter.
Why does this matter? What problem does it solve?]

%------------------------------------------------------------------------------
% INTUITION BOX: Motivating example
%------------------------------------------------------------------------------
\begin{intuition}[Why [Topic] Matters]
[Concrete example with named characters (Anna, Thomas, etc.)
that illustrates the core concept before any formalism.
Make the reader FEEL the problem before analyzing it.]
\end{intuition}

%------------------------------------------------------------------------------
% CENTRAL QUESTION BOX
%------------------------------------------------------------------------------
\paragraph{The Central Question.}
This chapter addresses the [N]th of the 9 BCM Questions:

\begin{center}
\fcolorbox{[COLOR]!75!black}{[COLOR]!5!white}{\parbox{0.85\textwidth}{
\textbf{Question [N]:} [Full question text from the 9 BCM Questions]
}}
\end{center}

%------------------------------------------------------------------------------
% CHAPTER OVERVIEW
%------------------------------------------------------------------------------
\paragraph{Chapter Overview.}
This chapter proceeds as follows:
\begin{itemize}
    \item Section \ref{sec:[label]-[topic1]}: [Description]
    \item Section \ref{sec:[label]-[topic2]}: [Description]
    \item Section \ref{sec:[label]-formal}: [Description]
    \item Section \ref{sec:[label]-applications}: [Description]
    \item Section \ref{sec:[label]-summary}: [Description]
\end{itemize}


%%%%%%%%%%%%%%%%%%%%%%%%%%%%%%%%%%%%%%%%%%%%%%%%%%%%%%%%%%%%%%%%%%%%%%%%%%%%%%%
%                     PART II: CORE CONTENT
%%%%%%%%%%%%%%%%%%%%%%%%%%%%%%%%%%%%%%%%%%%%%%%%%%%%%%%%%%%%%%%%%%%%%%%%%%%%%%%

\section{[First Major Topic]}
\label{sec:[label]-[topic1]}

%------------------------------------------------------------------------------

% Pattern for each subsection:
% 1. Definition environment for formal concept
% 2. Paragraph explaining characteristics
% 3. Paragraph with typical statement/example
% 4. Cross-reference to relevant appendix

\subsection{[Subtopic 1]}
\label{sec:[label]-[subtopic1]}

\begin{definition}[[Concept Name]]
[Formal definition in accessible language]
\begin{equation}
[Key equation]
\label{eq:[name]}
\end{equation}
\end{definition}

\paragraph{Characteristics.}
\begin{itemize}
    \item [Property 1]
    \item [Property 2]
    \item [Property 3]
\end{itemize}

\paragraph{Typical Statement.}
\textit{"[Quote that exemplifies this concept]"}

\paragraph{Intervention Focus.}
[What actions are appropriate at this point]

\begin{cross-reference}
See Appendix [CODE] ([NAME]) for [what formal content].
\end{cross-reference}

%------------------------------------------------------------------------------

\subsection{[Subtopic 2]}
\label{sec:[label]-[subtopic2]}

% Continue with same pattern...

\begin{intuition}[[Insight Title]]
[Explain WHY this works the way it does.
Use analogies, everyday examples, contrast with intuition.]
\end{intuition}


%%%%%%%%%%%%%%%%%%%%%%%%%%%%%%%%%%%%%%%%%%%%%%%%%%%%%%%%%%%%%%%%%%%%%%%%%%%%%%%
%                     PART III: FORMAL SPECIFICATION
%%%%%%%%%%%%%%%%%%%%%%%%%%%%%%%%%%%%%%%%%%%%%%%%%%%%%%%%%%%%%%%%%%%%%%%%%%%%%%%

\section{Formal Specification}
\label{sec:[label]-formal}

%------------------------------------------------------------------------------

% For TYPE A (CORE) chapters: Include key formal elements
% Keep derivations in appendix; show only essential equations here

\subsection{[Core Formalism 1]}

\begin{definition}[[Formal Concept]]
[Statement with equation]
\begin{equation}
[Equation]
\label{eq:[name]}
\end{equation}
where [define all symbols].
\end{definition}

\subsection{[Core Formalism 2]}

% Tables for structured parameter information
\begin{table}[h]
\centering
\begin{tabular}{lccc}
\toprule
\textbf{[Row Header]} & \textbf{[Col 1]} & \textbf{[Col 2]} & \textbf{[Col 3]} \\
\midrule
[Row 1] & [value] & [value] & [value] \\
[Row 2] & [value] & [value] & [value] \\
\bottomrule
\end{tabular}
\caption{[Table description]}
\label{tab:[name]}
\end{table}


%%%%%%%%%%%%%%%%%%%%%%%%%%%%%%%%%%%%%%%%%%%%%%%%%%%%%%%%%%%%%%%%%%%%%%%%%%%%%%%
%                     PART IV: INTEGRATION
%%%%%%%%%%%%%%%%%%%%%%%%%%%%%%%%%%%%%%%%%%%%%%%%%%%%%%%%%%%%%%%%%%%%%%%%%%%%%%%

\section{Integration with the CORE Framework}
\label{sec:[label]-integration}

%------------------------------------------------------------------------------

% For TYPE A chapters: Show how this CORE connects to all others

\subsection{Connection to [Earlier CORE]}
\label{sec:[label]-connection1}

\begin{cross-reference}
The [concept] from [CORE-NAME] determines:
\begin{itemize}
    \item [Connection 1]
    \item [Connection 2]
    \item [Connection 3]
\end{itemize}
See Chapter [N] and Appendix [CODE] for full specification.
\end{cross-reference}

\subsection{The 9C Integration Table}

% Show how this chapter's CORE relates to all 8 COREs
\begin{table}[h]
\centering
\begin{tabular}{lll}
\toprule
\textbf{CORE} & \textbf{Question} & \textbf{[This Chapter]'s Role} \\
\midrule
WHO (AAA) & Who has utility? & [How this chapter relates] \\
WHAT (C) & What is utility? & [How this chapter relates] \\
HOW (B) & How interact? & [How this chapter relates] \\
WHEN (V) & When context? & [How this chapter relates] \\
WHERE (BBB) & Where parameters? & [How this chapter relates] \\
AWARE (AU) & How aware? & [How this chapter relates] \\
READY (AV) & How willing? & [How this chapter relates] \\
STAGE (AW) & Where in journey? & [How this chapter relates] \\
\bottomrule
\end{tabular}
\caption{Integration of [this CORE] with the 9C framework}
\label{tab:8c-integration}
\end{table}


%%%%%%%%%%%%%%%%%%%%%%%%%%%%%%%%%%%%%%%%%%%%%%%%%%%%%%%%%%%%%%%%%%%%%%%%%%%%%%%
%                     PART V: WORKED EXAMPLES
%%%%%%%%%%%%%%%%%%%%%%%%%%%%%%%%%%%%%%%%%%%%%%%%%%%%%%%%%%%%%%%%%%%%%%%%%%%%%%%

\section{Worked Examples}
\label{sec:[label]-applications}

%------------------------------------------------------------------------------

% REQUIRED: At least one concrete example with numbers

\subsection{Example 1: [Named Character]'s [Situation]}

\paragraph{Baseline: [Initial State].}
\begin{itemize}
    \item Statement: "[Quote]"
    \item Parameters: $[param] = [value]$, $[param] = [value]$
    \item Analysis: [What this means]
\end{itemize}

\paragraph{Intervention: [What happens].}
\begin{itemize}
    \item "[Description of intervention]"
    \item Effect: [Parameter changes]
\end{itemize}

\paragraph{After Intervention: [New State].}
\begin{itemize}
    \item Statement: "[New quote]"
    \item Parameters: $[param] = [value]$, $[param] = [value]$
    \item Analysis: [What this means]
\end{itemize}

\begin{equation}
[Summary equation showing change]
\label{eq:[example-name]}
\end{equation}

%------------------------------------------------------------------------------

\subsection{Example 2: [Multi-Level / Different Domain]}

\begin{intuition}[[Title]]
[Show how the concept applies at different levels or domains]
\begin{itemize}
    \item $L=0$ (Individual): [Example]
    \item $L=1$ (Household): [Example]
    \item $L=2$ (Organization): [Example]
    \item $L=\infty$ (Society): [Example]
\end{itemize}
\end{intuition}

\begin{cross-reference}
Appendix [CODE] contains [additional examples/boundary conditions].
\end{cross-reference}


%%%%%%%%%%%%%%%%%%%%%%%%%%%%%%%%%%%%%%%%%%%%%%%%%%%%%%%%%%%%%%%%%%%%%%%%%%%%%%%
%                          PART VI: CLOSING
%%%%%%%%%%%%%%%%%%%%%%%%%%%%%%%%%%%%%%%%%%%%%%%%%%%%%%%%%%%%%%%%%%%%%%%%%%%%%%%

\section{Summary}
\label{sec:[label]-summary}

%------------------------------------------------------------------------------

This chapter has established:

\begin{enumerate}
    \item \textbf{[Key Point 1]:} [One sentence explanation]

    \item \textbf{[Key Point 2]:} [One sentence explanation]

    \item \textbf{[Key Point 3]:} [One sentence explanation]

    \item \textbf{[Key Point 4]:} [One sentence explanation]

    \item \textbf{[Key Point 5]:} [One sentence explanation]
\end{enumerate}

%------------------------------------------------------------------------------
\paragraph{Formal Details.}
The complete formal treatment appears in Appendix~[CODE] ([NAME]), which contains:
\begin{itemize}
    \item [Number] axioms specifying [what]
    \item [Number] theorems/boundary conditions
    \item [Specific formal content]
    \item [Novel predictions]
\end{itemize}

%------------------------------------------------------------------------------
\paragraph{What Comes Next.}
Chapter [N+1] introduces \textbf{[Next Chapter Title]}, which [brief description
of how it builds on this chapter].

%------------------------------------------------------------------------------
\begin{cross-reference}
\textbf{Reading Path:}
\begin{itemize}
    \item For formal [topic] specification: Appendix [CODE] ([NAME])
    \item For [related topic]: Chapter [N] + Appendix [CODE]
    \item For [related topic]: Chapter [N] + Appendix [CODE]
    \item For [next step]: Chapter [N+1] (next)
\end{itemize}
\end{cross-reference}

%%%%%%%%%%%%%%%%%%%%%%%%%%%%%%%%%%%%%%%%%%%%%%%%%%%%%%%%%%%%%%%%%%%%%%%%%%%%%%%
% END OF CHAPTER [N]
%%%%%%%%%%%%%%%%%%%%%%%%%%%%%%%%%%%%%%%%%%%%%%%%%%%%%%%%%%%%%%%%%%%%%%%%%%%%%%%


% *****************************************************************************
% *****************************************************************************
%
%                    STYLE GUIDELINES & REFERENCE
%
% *****************************************************************************
% *****************************************************************************


% =============================================================================
% PROSE STYLE
% =============================================================================
%
% - Write for intelligent non-specialists
% - Define technical terms when first used
% - Prefer active voice
% - Use short paragraphs (3-5 sentences)
% - Start with intuition, end with formalism
% - Use named characters (Anna, Thomas) in examples


% =============================================================================
% MATHEMATICAL STYLE
% =============================================================================
%
% - Number all important equations with \label{}
% - Provide interpretation after EVERY equation
% - Use consistent notation (see Appendix G: Glossary)
% - Keep derivations in appendices
% - Define symbols immediately after equations


% =============================================================================
% ENVIRONMENTS
% =============================================================================
%
% \begin{definition}[Name]    - Formal definitions, equations
% \begin{intuition}[Title]    - Conceptual explanations, WHY
% \begin{cross-reference}     - Links to other chapters/appendices
% \begin{tcolorbox}           - Highlighted boxes (CORE, Appendix refs)
%
% \paragraph{Title.}          - Key points within sections
% \textbf{term}               - First use of important terms
% \textit{quote}              - Typical statements, examples


% =============================================================================
% COLOR CODING FOR TCOLORBOX
% =============================================================================
%
% CORE Connection Boxes:
%   blue    - CORE-HOW (B)
%   green   - CORE-WHEN (V)
%   red     - CORE-WHAT (C)
%   purple  - CORE-WHO (AAA)
%   orange  - CORE-AWARE (AU)
%   violet  - CORE-READY (AV)
%   teal    - CORE-STAGE (AW)
%
% Standard Boxes:
%   yellow/orange - Appendix References
%   blue          - Key Insights
%   gray          - Examples (mdframed)


% =============================================================================
% CROSS-REFERENCES
% =============================================================================
%
% - Always specify appendix by BOTH code AND name: "Appendix C (CORE-WHAT)"
% - Reference chapters by number: "as shown in Chapter 5"
% - Use \ref{} for sections: "Section \ref{sec:label}"
% - Use \eqref{} for equations: "Equation \eqref{eq:name}"
% - Use \pageref{} sparingly for distant references


% =============================================================================
% LENGTH GUIDELINES
% =============================================================================
%
% - Header (metadata, boxes): 1-2 pages
% - Opening section: 1-2 pages
% - Each content section: 3-8 pages
% - Worked examples: 2-4 pages
% - Summary: 1-2 pages
% - Total chapter: 15-50 pages (varies by complexity)


% =============================================================================
% CHAPTER-APPENDIX RELATIONSHIP
% =============================================================================
%
% WHAT GOES IN CHAPTERS:
% - Conceptual motivation
% - Intuitive explanations
% - Simple examples with named characters
% - Key equations (with interpretation)
% - Policy implications
% - Connections to literature (brief)
%
% WHAT GOES IN APPENDICES:
% - Complete axiom systems (8-200 axioms)
% - Full proofs and derivations
% - All edge cases and boundary conditions
% - Comprehensive literature review (80-150+ refs)
% - Detailed worked examples
% - Symbol glossaries
% - Critical foundations (objections/responses)
% - Open issues and research roadmap
%
% REFERENCE PATTERN:
% In chapter: "The utility function has six dimensions (see Appendix C for
%              the complete 144-component specification)."
% In appendix: "This appendix provides the formal details for Section 10.2
%               of the main text."


% =============================================================================
% EBF CHAPTER LIST WITH CORE MAPPINGS
% =============================================================================
%
% FOUNDATION CHAPTERS (Type B):
% Ch. 1:  Introduction
% Ch. 2:  Rationality as Stability
% Ch. 3:  Limits of Classical Utility
% Ch. 4:  Empirical Foundations
% Ch. 4x: Calibration (LLM Monte Carlo)         → METHOD-LLMMC (AN)
% Ch. 6:  Reference Structure
% Ch. 7:  Fit and Non-Concavity
% Ch. 8:  Mathematical Foundations              → FORMAL-DERIVE (A)
%
% CORE CHAPTERS (Type A):
% Ch. 5:  Complementarity                       → CORE-HOW (B)
% Ch. 9:  Context                               → CORE-WHEN (V)
% Ch. 10: Welfare (FEPSDE)                      → CORE-WHAT (C), CORE-WHO (AAA)
% Ch. 11: Awareness                             → CORE-AWARE (AU)
% Ch. 12: Willingness                           → CORE-READY (AV)
% Ch. 13: Behavioral Change Journey             → CORE-STAGE (AW)
%
% APPLICATION CHAPTERS (Type C):
% Ch. 14: Behavioral Change Segments
% Ch. 15: WEC Synthesis
% Ch. 16: Probability of Change
% Ch. 17: Policy Implications
% Ch. 18: Limitations
% Ch. 19: Conclusion
%
% =============================================================================


% =============================================================================
% CHECKLIST: BEFORE FINALIZING A CHAPTER
% =============================================================================
%
% HEADER:
% [ ] Metadata block complete (version, purpose, appendices, prerequisites)
% [ ] \chapter{} and \label{} present
% [ ] \BCMQuestionNav{} with correct chapter number
% [ ] CORE Connection Box (if Type A chapter)
% [ ] Appendix References Box with section links
%
% OPENING:
% [ ] 2-3 sentence opening paragraph
% [ ] Intuition box with named character example
% [ ] Central Question box
% [ ] Chapter Overview with section references
%
% CONTENT:
% [ ] All subsections have \label{}
% [ ] Definition environments for formal concepts
% [ ] Intuition environments for explanations
% [ ] Cross-reference environments for links
% [ ] All equations numbered and interpreted
% [ ] Tables properly formatted with \caption and \label
%
% EXAMPLES:
% [ ] At least one worked example with numbers
% [ ] Named characters (Anna, Thomas, etc.)
% [ ] Before/after comparison
%
% CLOSING:
% [ ] Summary with 4-6 enumerated key points
% [ ] Formal Details paragraph with appendix content list
% [ ] What Comes Next paragraph with transition
% [ ] Reading Path cross-reference box
%
% STYLE:
% [ ] Consistent notation (checked against Appendix G)
% [ ] All appendix references include code AND name
% [ ] No orphaned cross-references (\ref{} to non-existent labels)
% [ ] Chapter length within 15-50 pages
%
% =============================================================================
